\section{Details Omitted from the Main Text}
\label{app:details}
\emph{In short: the complete text of the paragraphs that are condensed in the main text, in the order of the sections.}

\subsection{Introduction: scope, channels and roadmap}
In this work we use the LHCb Open Data record for $B^{\pm}\to K^{\pm}K^{+}K^{-}$~\cite{opendata_record,opendata_notebook}, a subset of the 2011 data on which the first evidence of an inclusive $CP$ asymmetry in this decay was published~\cite{lhcb_cpv_2013}, to address the question: \emph{to what extent do frequentist likelihood fits and Bayesian inference produce consistent, well-controlled estimates of the raw asymmetry $\Araw$ in this decay?} The aim is not a competitive measurement of $\ACP$, but a methodological comparison of the two inference frameworks on the same data and the same physical question. Of the four channels in the record, $B^\pm\to K^\pm K^+K^-$ has the lowest background; the others ($\pi^\pm\pi^+\pi^-$, $K^\pm\pi^+\pi^-$, $\pi^\pm K^+K^-$) need a different particle-identification selection and a model of larger and peaking (misidentified) backgrounds, and are not analysed, to keep the focus on the comparison of the methods. The published $\ACP$ cannot be reproduced from this record, which does not contain the control channel used to correct $\Araw$ for the production and detection asymmetries (Appendix~\ref{sec:limitations}); what we extend is the statistical treatment: a common model fitted with two inference methods, a common systematic study, a validation of the intervals, and a machine-learning (ML) selection, described in Appendix~\ref{sec:ml}, whose effect on the measurement is tested. Section~\ref{sec:data} describes the dataset and the preselection common to the analyses and Sec.~\ref{sec:model} the mass model and likelihood common to the two mass fits; Secs.~\ref{sec:frequentist} and \ref{sec:bayesian} present the frequentist and Bayesian fits, Sec.~\ref{sec:syst} their common systematic study, Sec.~\ref{sec:validation} the validation of the estimator and of the intervals with toy-MC experiments, Sec.~\ref{sec:dalitz} the local asymmetry in the Dalitz plane, and Sec.~\ref{sec:comparison} compares the two approaches; the machine-learning selection is in Appendix~\ref{sec:ml}. The main results are collected in Table~\ref{tab:summary}.

\subsection{Extended likelihood versus a mixture with fixed total}
The extended form differs from a mixture of signal and background with a fixed total number of candidates and a signal fraction only through the Poisson term for the total yield, which carries no information on $A$ and does not change its estimate; it is used because the yields of the two charges are the natural parameters of the model. Bayes' theorem applied to the fitted yields gives the probability that a candidate in the signal window $|m-\mu|<2\sigma$ is a signal decay, about $96\%$: a purity, not a property of the single candidate, which depends on the signal fraction of the sample (here $78\%$ of the candidates in the fit window).

\subsection{MCMC sampling and convergence (complete)}
The posterior was sampled with the affine-invariant ensemble sampler of Goodman and Weare~\cite{goodman_weare} as implemented in \texttt{emcee}~\cite{emcee} (the stretch move, the autocorrelation time and the effective sample size are described in Appendix~\ref{app:theory}). The maximum of the posterior was first located numerically (Nelder-Mead), and 32 walkers were initialized in a narrow ball around it and run for $25{,}000$ steps, with the first $4166$ ($1/6$) discarded as burn-in, leaving $6.7\times10^5$ samples; the mean acceptance fraction is $0.38$. Convergence was checked through the integrated autocorrelation time of each parameter, $\tau=1+2\sum_{t\geq1}\rho(t)$, with $\rho(t)$ the autocorrelation of the chain at lag $t$. We find $\tau\in[147,183]$ steps, the longest for the background-shape parameters: the post-burn-in chain is $N/\tau\geq113$ autocorrelation times long for every parameter, above the threshold of $50$ recommended for a reliable estimate of $\tau$~\cite{emcee}, and corresponds to $\gtrsim3600$ effectively independent samples per parameter. To check that the chain is not trapped near its starting point, the fit was repeated twice with the walkers initialized uniformly over the whole prior box and the first half of the chain discarded: both runs converge to the same single mode, with $\Araw=-0.0446$ and $-0.0442$, the same intervals and the same values of all other parameters as the nominal fit. The trace plot of such a run (Fig.~\ref{fig:bayes_trace}, Appendix~\ref{sec:appendix}) shows all walkers reaching the bulk of the posterior within $\sim$2000 steps, the background fraction $f_+$ last, well inside the discarded half of the chain.

\subsection{Posterior of the raw asymmetry (complete)}
The marginal posterior on the raw asymmetry gives (median and equal-tailed credible intervals, CI)
\begin{gather*}
\Araw = -0.0443, \\
68\%~\text{CI}~[-0.0528,-0.0358], \\ 95\%~\text{CI}~[-0.0608,-0.0276].
\end{gather*}
The posterior is very close to Gaussian, with standard deviation $0.0085$; its 68\% interval is symmetric about the median to within $2\times10^{-5}$, so the equal-tailed intervals coincide with the highest-posterior-density ones at the quoted precision. Since the priors are uniform, the posterior mode (the maximum a posteriori, MAP, estimate, $-0.0441$) coincides with the maximum-likelihood estimate, so this result can be compared directly with the frequentist fit of the same model (Sec.~\ref{sec:frequentist}; the quantities it uses are summarized in Appendix~\ref{sec:appendix_freq}). The published analysis of the same 2011 data, with its own selection (which includes a $D^0$ veto) and a mass model with partially reconstructed and misidentified backgrounds, finds $\Araw=-0.060\pm0.007$~\cite{lhcb_cpv_2013}; with the $D^0$ veto our fit gives $-0.0496\pm0.0088$ (Sec.~\ref{sec:bayes_checks}), compatible within the statistical uncertainties, although the two selections overlap and differ. After the corrections of Eq.~\eqref{eq:acp}, the published asymmetry is $\ACP=-0.043\pm0.009\pm0.003\pm0.007$~\cite{lhcb_cpv_2013}, and the most precise later measurement, on a larger data set, is $\ACP=-0.037\pm0.002\pm0.002\pm0.003$~\cite{lhcb_cpv_prd2023,lhcb_cpv_paper}. The other physics parameters are $N_s=15{,}328\pm133$, $N_b=4404\pm82$, $A_b=-0.040\pm0.018$, $\mu=5284.39\pm0.16\,\MeV$ and $\sigma=18.00\pm0.14\,\MeV$ (Fig.~\ref{fig:bayes_corner}). $A$ is almost uncorrelated with the other parameters: its largest correlation, $-0.17$, is with $A_b$: conditioning on $A_b=0$ shifts $A$ by $\rho\,\sigma_A/\sigma_{A_b}\,(0-\bar A_b)=-0.0033$, to $-0.0475$, which is what the fit with a charge-symmetric background gives ($-0.0471$, Sec.~\ref{sec:syst}).

\subsection{Bayes factor of A=0 (complete)}
The Bayes factor compares two hypotheses through their evidences (marginal likelihoods), $BF_{10}=P(D\mid H_1)/P(D\mid H_0)$ with $P(D\mid H_i)=\int\mathcal{L}(D\mid\theta,H_i)\,\pi(\theta\mid H_i)\,d\theta$: it is the factor by which the data update the prior odds of the two hypotheses, and it requires proper priors. Computing the evidence of each hypothesis by direct numerical integration is impractical for a 12-parameter posterior. When the null hypothesis fixes one parameter, $H_0$: $A=0$, nested in $H_1$ ($A$ free), and the priors on the other parameters are shared, the Bayes factor reduces to the Savage-Dickey density ratio~\cite{dickey1971} (Appendix~\ref{sec:appendix_freq}),
\[
BF_{01} = \frac{p(A=0\mid D)}{\pi(A=0)}, \qquad \pi(A=0)=\tfrac12 .
\]
With equal prior probabilities for the two hypotheses, $BF_{01}$ is also the posterior odds of $H_0$ against $H_1$. $A=0$ lies $5.2$ posterior standard deviations from the median, in a tail where the chain contains essentially no samples, so a kernel density estimate of $p(A=0\mid D)$ would only reflect the kernel tails of the few most extreme samples. We therefore use the Gaussian approximation of the marginal posterior (Sec.~\ref{sec:bayes_posterior}), with the mean and standard deviation of the chain, which gives $BF_{01}\approx1.1\times10^{-4}$ ($BF_{10}\approx9\times10^{3}$, $\ln BF_{10}\approx9.1$): decisive evidence against $A=0$ on Jeffreys' scale~\cite{jeffreys} ($BF_{10}>100$). As for any Bayes factor of a point null against a free parameter, the value scales with the prior width -- the Ockham factor built into the Bayes factor, which penalizes a free parameter in proportion to the prior volume it is allowed to explore, and the mechanism behind Lindley's paradox: with the narrower prior $A\in(-0.3,0.3)$ of Sec.~\ref{sec:bayes_checks}, $\pi(A=0)$ grows to $1/0.6$ and $BF_{01}$ drops by a factor $3.3$, which leaves the conclusion unchanged. Figure~\ref{fig:bayes_sd} of Appendix~\ref{app:ext_bayes} shows the two densities of the ratio; the dependence on the prior width, the evidence computed by direct integration and the equality of the background shapes of the two charges are in Appendix~\ref{sec:appendix_checks}.

\subsection{Bayes factor of the background asymmetry (complete)}
The same test on the background asymmetry shows the Ockham factor at work. $A_b=-0.040\pm0.018$ lies $2.3$ standard deviations from zero, yet the Savage-Dickey ratio gives $BF_{01}(A_b=0)\approx3.5$: the data mildly \emph{favor} a charge-symmetric background, because the gain in likelihood does not compensate for spreading the prior of $A_b$ over $(-1,1)$. A negative background asymmetry is plausible: besides combinatorial background, the sample contains misidentified $B^\pm\to K^\pm\pi^+\pi^-$ and $B^\pm\to\pi^\pm K^+K^-$ decays, with $\ACP$ values of their own, whose contribution grows when the kaon identification is loosened (Sec.~\ref{sec:syst}). $A_b$ is nevertheless kept free in the nominal model, which avoids imposing an assumption on the background that the frequentist analysis does not make; fixing $A_b=0$ is one of the variants of the systematic study (Sec.~\ref{sec:syst}).

\subsection{Combination rule of the systematic study (complete)}
The shift $\Delta=\Araw^{\mathrm{var}}-\Araw^{\mathrm{nom}}$ of a variant is not necessarily a systematic effect. When a variant changes the sample (window, preselection), the two samples are nested and part of the shift is a statistical fluctuation, with expected size $\sigma_B=\sqrt{|\sigma_{\mathrm{var}}^2-\sigma_{\mathrm{nom}}^2|}$, $\sigma$ being the statistical uncertainty of each fit~\cite{barlow}; such a variant contributes only if $|\Delta|>\sigma_B$. A variant that the data reject, with a $\chi^2$ worse than the nominal one by more than $25$ on the same data, is not a plausible alternative and is excluded: this removes the single exponential ($\Delta\chi^2=134$, Sec.~\ref{sec:bayes_gof}). The largest remaining $|\Delta|$ of each group is taken as its uncertainty, since the variants of a group are alternative descriptions of the same choice, and the four groups are added in quadrature.

\subsection{Systematic study: results and the particle-identification scan (complete)}
The results are in Table~\ref{tab:syst}, and all the Bayesian fits, with their intervals and goodness of fit, in Table~\ref{tab:bayes_all} of the Appendix. The two methods give the same shifts to within $\sim$$0.0005$ and the same group uncertainties, with total systematic uncertainties of $0.0070$ (Bayesian) and $0.0072$ (frequentist). The final results are
\begin{gather*}
\text{Bayesian:}\\ \Araw=-0.0443\pm0.0085\,(\text{stat})\pm0.0070\,(\text{syst}),\\[2pt]
\text{frequentist:}\\ \Araw=-0.0441\pm0.0085\,(\text{stat})\pm0.0072\,(\text{syst}).
\end{gather*}
The preselection dominates, through the single variant $\texttt{ProbK}>0.6$ ($+0.0058$, $1.2\,\sigma_B$). To tell a dependence on the PID requirement from a fluctuation of that single value, the frequentist fit was repeated in steps of $0.05$ of each cut, from $\texttt{ProbK}>0.30$ to $>0.85$ and from $\texttt{ProbPi}<0.15$ to $<0.70$. On the tight side of the kaon requirement the shifts scatter around zero without a trend ($+0.0009$, $+0.0058$, $-0.0017$, $-0.0112$, $-0.0054$, $+0.0010$, $+0.0022$ for $\texttt{ProbK}>0.55$--$0.85$, all within $1.3\,\sigma_B$), and the pion requirement gives a flat scan within $1.3\,\sigma_B$: the $\texttt{ProbK}>0.6$ value is a statistical fluctuation, and the preselection uncertainty is a conservative estimate. Loosening the kaon requirement, instead, moves $\Araw$ monotonically ($-0.0019$, $-0.0032$, $-0.0067$, $-0.0111$ for $\texttt{ProbK}>0.45$--$0.30$, the last at $2.5\,\sigma_B$), as expected if misidentified $B\to K\pi\pi$ and $B\to\pi KK$ decays, with asymmetries of their own and not included in the mass model, enter the sample; this supports the nominal requirement. The combination rule was fixed before this scan and is not changed by it. The background group is the second contribution and is driven by the assumptions on the charge dependence of the background, the window the third.

\subsection{Sensitivity of sigma B to small differences between the methods}
Each method computes $\sigma_B$ from its own uncertainties (posterior standard deviation, Hesse error). Since $\sigma_B$ is the difference of two nearly equal uncertainties, it is sensitive to small differences between them, including the Monte Carlo noise of the posterior standard deviation (1--3\%), and a variant close to the threshold can be classified differently by the two methods: $\texttt{ProbPi}<0.4$ counts for the Bayesian fit ($|\Delta|=0.00223$, $\sigma_B=0.00220$) and not for the frequentist one ($0.00204$ vs $0.00229$). This does not affect the result, since the uncertainty of each group is set by its largest shift, the same in both methods.

\subsection{Stability scans of the frequentist fit: numbers (complete)}
\paragraph*{Numbers.} Three results are used here, from the same notebook as the nominal fit; they do not enter the systematic uncertainty, whose variants and combination rule are those above.
\begin{itemize}[leftmargin=*,itemsep=1pt]
\item \emph{Fit window}: the maximum shift with respect to the nominal fit is $0.0048$ over the grid of lower and upper edges (standard deviation $0.0018$, mean statistical uncertainty $0.0085$); the extreme corners are the smallest and most asymmetric windows, where the background shape is least constrained.
\item \emph{Signal region} (used only to start the sideband fit): the shift is at most $6\times10^{-5}$, so the result does not depend on it.
\item \emph{Particle-identification requirements}: on the grid of the two cuts the maximum shift is $0.0126$ (standard deviation $0.0046$). This is larger than the statistical uncertainty of the nominal fit, $0.0085$, but the grid contains much smaller samples, with a mean statistical uncertainty of $0.0112$; the relevant comparison for nested samples is with $\sigma_B$, as in the one-cut-at-a-time scan above, where no shift exceeds $1.3\,\sigma_B$ on the tight side.
\end{itemize}
The window and signal-region scans are therefore consistent with a fit that does not depend on these choices, and the preselection remains the dominant contribution, as quantified by the systematic study.

\subsection{Validation: generation of the toy MC experiments (complete)}
The frequentist results of Sec.~\ref{sec:frequentist} rest on asymptotic properties of the maximum-likelihood estimator and of the profile-likelihood ratio, and Sec.~\ref{sec:freq_coverage} verified the coverage by resampling the data. Here the same questions are asked of toy-MC experiments generated from the fitted model, as was done for the Bayesian interval in Sec.~\ref{sec:bayes_toymc}, which allows the true value of $A$ to be chosen. The fitted model of Sec.~\ref{sec:model} defines the toy-MC datasets: for each charge the yields of signal and background are Poisson-fluctuated, the signal masses are drawn from the truncated Gaussian and the background masses from the two truncated exponentials by the transformation method. Each dataset is fitted with the nominal frequentist procedure; the likelihood ratio $t(A_{\mathrm{true}})=-2\ln[L(A_{\mathrm{true}})/L(\hat A)]$ and the statistic of the test of $A=0$, $q(0)$, are computed with the other eleven parameters profiled. One thousand toy-MC experiments were generated at the fitted value, $A_{\mathrm{true}}=-0.0441$, and $300$ at each of $A_{\mathrm{true}}=-0.1$, $-0.02$, $0$ and $+0.05$; between $1\%$ and $3\%$ of the fits are not flagged as valid by \texttt{MIGRAD} and are excluded.

\subsection{Validation: bias, efficiency and coverage (complete)}
Table~\ref{tab:validation_toys} summarizes the estimator over the five injected values. The bias of $\hat A$ is at most $0.0008$ ($1.6$ standard errors), the spread of the estimates equals the mean Hesse uncertainty within $6\%$ (ratio $0.94$--$1.05$), and the pulls $(\hat A-A_{\mathrm{true}})/\sigma(\hat A)$ have width $0.94$--$1.05$: the estimator is unbiased and reaches the minimum variance bound for this sample size, and the Hesse uncertainty is the correct statistical uncertainty. At the fitted value the coverage of the interval $t(A_{\mathrm{true}})\leq1$ is $0.696\pm0.015$ for the $68.3\%$ level and $0.954\pm0.007$ for the $95.4\%$ level (Fig.~\ref{fig:validation_null} of Appendix~\ref{app:ext_valid}, panel b, shows the $p$-value of $t(A_{\mathrm{true}})$, uniform as expected, Kolmogorov-Smirnov $p=0.96$ for all values pooled). The $68\%$ coverage at $A_{\mathrm{true}}=+0.05$, $0.614\pm0.028$, is $2.5$ standard deviations below nominal; the $95\%$ coverage at the same value is $0.956$, and with five injected values a deviation of this size is not unexpected.

\subsection{Dalitz plane: local deviations and their robustness (complete)}
Two of the 32 bins have $|S_i|>3$, adjacent to each other at $m^2_{\mathrm{lo}}<4\,\GeV^2$ and $7.5<m^2_{\mathrm{hi}}<11.25\,\GeV^2$ ($S_i=-3.7$ and $-3.0$), with $A_i$ more negative than the global value; the bin below the first one ($m^2_{\mathrm{lo}}<2$, $3.75<m^2_{\mathrm{hi}}<7.5\,\GeV^2$) has $S_i=-2.98$, just under the threshold (Fig.~\ref{fig:bayes_dalitz}); about $0.09$ would be expected if the asymmetry were uniform, with a probability of $3\times10^{-3}$ for two or more. Because the bins are examined together this is a look-elsewhere situation: the number refers to any two of the $32$ bins, not to the two found, and the global homogeneity test that follows does not depend on which bins are singled out. The largest positive deviations, close to $+3$, lie at high $m^2_{\mathrm{hi}}$. A global test of homogeneity, $\chi^2=\sum_i S_i^2=80.8$ for 31 degrees of freedom ($p=3\times10^{-6}$), confirms that the local asymmetry is not uniform. Restricting the sample to the signal window ($|m-\mu|<2\sigma$, $n=15{,}130$, where the fitted model gives a signal purity of $\approx$96\%) gives the same picture -- three bins with $|S_i|>3$, $\chi^2=95.6/28$ -- so the structure is not produced by the combinatorial background. Nor is it produced by the charm contributions: with the $D^0$ veto two bins still exceed $3\sigma$ ($\chi^2=79.2/31$), and with all three vetoes one remains ($\chi^2=74.4/31$, $p=2\times10^{-5}$). It does not depend on the binning either: $6\times6$, $10\times10$ and $12\times12$ grids give $\chi^2=81.0/19$, $97.1/43$ and $120.0/60$ ($p\leq7\times10^{-6}$).

\subsection{Dalitz plane: the region of the published analysis (complete)}
The ordered Dalitz plot (Fig.~\ref{fig:dalitz_region}a) shows the two-body resonances of the decay: the narrow $\phi(1020)$ band at the kinematic edge, $m^2_{\mathrm{lo}}\approx1.04\,\GeV^2$, and the charm bands of the $D^0$, $J/\psi$ and $\chi_{c0}$. The negative local asymmetry is not produced by the $\phi(1020)$: its band ($|m_{\mathrm{lo}}-m_\phi|<10\,\MeV$, $1964$ candidates) has $A_{\mathrm{count}}=-0.004\pm0.023$, while the significant bins without it have $-0.145\pm0.019$. The asymmetry sits in the low-mass $K^+K^-$ region just above the $\phi$, which is where the published analysis of the same data located it, defining the region $1.2<m^2_{\mathrm{lo}}<2.0\,\GeV^2$, $m^2_{\mathrm{hi}}<15\,\GeV^2$ and measuring there $\Araw=-0.239\pm0.020$~\cite{lhcb_cpv_2013}. Since that region is fixed by the paper and not chosen on our data, it allows an unbiased comparison: fitting the $2146$ candidates in it with the nominal model gives $\Araw=-0.234\pm0.025$ (Bayesian; the chain mixes slowly in the poorly constrained second background exponential, $N/\tau=15$) and $-0.237\pm0.025$ (frequentist, $9.4\sigma$ from zero), reproducing the published value; the mass spectra of the two charges in the region (Fig.~\ref{fig:dalitz_region}b) show the difference directly. Three caveats apply: the binning is coarse and was not optimized; the asymmetries are computed on all candidates, without the sWeight-based background subtraction that a full analysis would use (Appendix~\ref{sec:limitations}); and, since $\Araw$ is not corrected for detection asymmetries, local variations can still be partly instrumental -- the $K^+$/$K^-$ interaction asymmetry depends on the kaon momenta, which are correlated with the position in the Dalitz plane, and is not removed by the polarity split. Appendix~\ref{sec:ml_kmeans} revisits the same question with an adaptive, data-driven partition ($K$-means).

\subsection{Comparison: hypothesis test (complete)}
\emph{Hypothesis test.} The simplest test, the exact binomial test of $N^-=N^+$ on all candidates of the window, gives $6.1\sigma$ ($p=1.3\times10^{-9}$, $A_{\mathrm{count}}=-0.0432\pm0.0071$); it sees the asymmetry of signal and background together and does not use the mass. The frequentist test of the fit gives a $p$-value, $2\times10^{-7}$ ($Z=\sqrt{q(0)}=5.2\sigma$ in the two-sided convention; the one-sided $Z=\Phi^{-1}(1-p)$ would give $5.06\sigma$, Appendix~\ref{sec:appendix_freq}), which does not depend on the prior and is not the probability of $A=0$; its size and power were verified with toy-MC experiments (Sec.~\ref{sec:validation_power}). In particle physics a significance of $5\sigma$ is the conventional threshold for a discovery claim, chosen to limit the effect of multiple testing and of unforeseen systematic effects; the $5.2\sigma$ found here crosses it, but it refers to a raw asymmetry that includes production and detection effects (Sec.~\ref{sec:comparison_published}), and is not a discovery of CP violation. The Bayes factor, $BF_{10}\approx9\times10^3$, compares the evidence of $A$ free and $A=0$ and depends on the prior width: with $A\in(-0.3,0.3)$ instead of $(-1,1)$ it decreases by a factor $3.3$. For $A$ the two give the same conclusion. For the background asymmetry, $A_b=-0.040\pm0.018$, they would lead to different statements: a $2.3\sigma$ deviation from zero, which a frequentist test would call a hint, and $BF_{01}\approx3.5$, a mild preference for $A_b=0$ because a free parameter with a prior of width two is penalized (Sec.~\ref{sec:bayes_factor}). Neither is wrong; they answer different questions.

\subsection{Comparison: goodness of fit (complete)}
\emph{Goodness of fit.} The same fit of the same data gives $\chi^2/\mathrm{ndof}=441.7/388$ ($p=0.031$) in $2\times200$ bins of $2.5\,\MeV$ and $\approx160/88$ ($p\approx5\times10^{-6}$) in $2\times50$ bins of $10\,\MeV$ (the Bayesian fit, $164.1/88$, $p=1.6\times10^{-6}$, uses the second binning). The misfit is at the scale of $\sim10\,\MeV$, the flanks of the peak, and is diluted in narrow bins where Poisson fluctuations dominate; the $p$-value of a binned $\chi^2$ is a property of the binning as much as of the model. The Bayesian analysis adds the posterior predictive check, which accounts for the posterior uncertainty of the parameters and shows that the misfit is concentrated in the $B^+$ sample (Sec.~\ref{sec:bayes_gof}).

\subsection{Comparison with the published analysis: samples and models (complete)}
The published analysis of the same data~\cite{lhcb_cpv_2013} finds $\Araw=-0.060\pm0.007$ with $22{,}119\pm164$ signal candidates, against $-0.0441\pm0.0085$ and $15{,}328\pm133$ here. The two differ in the sample (the Open Data record is a preselected subset, reduced to $23{,}895$ candidates by the identification requirement alone) and in the model, which includes a Cruijff signal, partially reconstructed four-body decays and a misidentified peaking background from simulation. With the same $D^0$ veto as the published selection the Bayesian fit gives $-0.0496\pm0.0088$, within one standard deviation of the quadrature sum of the two uncertainties ($0.93\sigma$), although the samples overlap and the selections differ.

\subsection{Comparison with the published analysis: raw versus CP asymmetry (complete)}
The raw asymmetry is not the CP asymmetry. The two are related by $\ACP\approx\Araw-A_P-A_D$ (Eq.~\eqref{eq:acp}) and, in the published analysis, the correction obtained from $B^\pm\to J/\psi K^\pm$ decays and from the acceptance moves the result from $-0.060$ to $-0.043\pm0.009\pm0.003\pm0.007$ ($3.7\sigma$)~\cite{lhcb_cpv_2013}, i.e.\ by $+0.017$, of the order of two statistical uncertainties of this work. A correction of that size cannot be assumed for a different selection, and the $5.2\sigma$ of this work is the significance of a nonzero raw asymmetry, which includes the production asymmetry of $B^\pm$ mesons and the different interaction of $K^+$ and $K^-$ with the detector, not a significance of CP violation. A later measurement on a larger data set gives $\ACP=-0.037\pm0.002\pm0.002\pm0.003$ ($8.5\sigma$)~\cite{lhcb_cpv_prd2023,lhcb_cpv_paper}. Our result is consistent in sign and size with these values, which is what can be said without the control channel, absent from this record.

\subsection{Conclusions: headline (complete)}
We have measured the raw charge asymmetry of $B^\pm\to K^\pm K^+K^-$ decays in the 2011 LHCb Open Data record with a common selection and a common mass model, and compared how two inference frameworks treat the same question. The frequentist profile-likelihood fit and the Bayesian posterior give $\Araw=-0.0441\pm0.0085\pm0.0072$ and $-0.0443\pm0.0085\pm0.0070$ (statistical, systematic), a nonzero asymmetry at $5.2$ standard deviations that is stable under the thirteen variants of the fit and selection. The result is compatible in sign and size with the published analysis of the same data and with its later, more precise measurements, and an asymmetry concentrated in a region of the Dalitz plane, outside the $\phi(1020)$, is reproduced a priori in the region defined by the published analysis. It is a raw asymmetry: without the control channel that corrects for the production of $B^\pm$ mesons and for the detection of $K^\pm$, it is not a measurement of $\ACP$, and its significance is not that of CP violation (Sec.~\ref{sec:comparison_published}).

\subsection{Conclusions: lessons of each method (complete)}
Each method taught something specific.
\begin{itemize}[leftmargin=*,itemsep=2pt]
\item \emph{Frequentist.} With the other eleven parameters profiled, Wilks' theorem gives a $68\%$ interval with the nominal coverage ($68.4\pm1.5\%$ in $1000$ bootstrap replicas of the data, $0.70\pm0.01$ and $0.95\pm0.01$ in toy-MC experiments from the fitted model, where the estimator is also unbiased and reaches the minimum variance bound), and the Hesse and Minos uncertainties agree, so the likelihood is parabolic in $A$ at this sample size.
\item \emph{Bayesian.} With $\sim2\times10^4$ candidates and flat priors the posterior is the likelihood, and the credible interval coincides with the confidence interval and has the same coverage in toy experiments. The Bayesian analysis adds what the frequentist one lacks: sampling diagnostics that must be checked rather than assumed (autocorrelation time, dispersed starts, label switching of the two background slopes), a Bayes factor that depends on the prior width, and a posterior predictive check whose charge statistics, being fitted, have no power.
\item \emph{Machine learning (Appendix~\ref{sec:ml}).} Weak labels from the mass spectrum are enough to train classifiers of test AUC $\simeq0.86$, but the selection does not improve the precision on $\Araw$ in a sample with little background, and it cannot remove the misidentified decays that peak under the signal. Its value is as an independent test of the stability of $\Araw$ under a different selection and, with $K$-means, as an adaptive description of the Dalitz plane that agrees with the fixed binning.
\end{itemize}

\subsection{Conclusions: systematic study and dependence on choices (complete)}
On the systematic study, the dominant contribution is the particle-identification requirement, which a finer scan shows to be consistent with a statistical fluctuation of a nested sample; the background model is second, and the signal shape does not matter. The same variants and the same combination rule give the same uncertainty in the two methods. More generally, several statistics that look informative depend on choices: the $p$-value of a binned $\chi^2$ on the binning, the Bayes factor on the prior, the coverage on whether the toy-MC experiments come from the model or from the data.

\subsection{Data: description of the record (complete)}
The dataset is the CERN Open Data record 4900, ``Matter Antimatter Differences -- B meson decays to three hadrons''~\cite{opendata_record}: candidates for $B^{\pm}\to h^{\pm}h^{+}h^{-}$ decays collected by LHCb in 2011 at $\sqrt{s}=7$\,TeV, reconstructed with the LHCb software and already passed through trigger and offline selection. The data are distributed as ROOT \texttt{TTree}s (\texttt{DecayTree}) containing only the most relevant reconstructed branches, split across two files by magnet polarity (magnet up/down), for a combined $8{,}556{,}118$ candidates. The same record also includes a simplified, detector-response-free phase-space simulation of $B^{\pm}\to K^{\pm}K^{+}K^{-}$ ($50{,}000$ decays, half of each charge); it is used only as a sample without asymmetry to test the local-asymmetry procedure (Sec.~\ref{sec:bayes_dalitz}).

\subsection{Data: branches and preselection (complete)}
For each of the three final-state tracks ($H1$, $H2$, $H3$) the analysis reads momentum components (\texttt{PX}, \texttt{PY}, \texttt{PZ}), particle-identification (PID) likelihoods (\texttt{ProbK}, \texttt{ProbPi}), a muon flag (\texttt{isMuon}), and the track charge. The same preselection is used in all three analyses: a kaon-PID requirement and a muon veto on every track,
\[
\texttt{H*\_ProbPi} < 0.5 \ \wedge\ \texttt{H*\_ProbK} > 0.5 \ \wedge\ \lnot\texttt{H*\_isMuon},
\]
which retains $23{,}895$ of the $8{,}556{,}118$ candidates ($0.28\%$), the PID requirement being by far the dominant reduction (Table~\ref{tab:cutflow} of Appendix~\ref{app:ext_data}). The files contain no event identifier; as a proxy for multiple candidates from the same collision, only $42$ of the candidates in the fit window ($0.2\%$) share an identical track with another candidate, a negligible effect on the statistical uncertainties, and all are kept.

\subsection{Priors: flat priors and the evidence (complete)}
A flat prior is not uninformative in an absolute sense, since it is not invariant under reparametrization (a prior uniform in $\lambda$ is not uniform in $1/\lambda$); with $\sim$$2\times10^4$ events, however, the posterior is dominated by the likelihood, as the prior-sensitivity check of Sec.~\ref{sec:bayes_checks} confirms for $A$. The evidence $p(D)$ is a constant normalization: the MCMC samples $\mathcal{L}\,\pi$ without computing it, and the Savage-Dickey ratio of Sec.~\ref{sec:bayes_factor} avoids computing it as well. The marginal posterior of the asymmetry, $p(A\mid D)=\int p(\theta\mid D)\,d\theta_{\neq A}$, is the histogram of the $A$ component of the chain.

\subsection{Comparison: meaning of the interval (complete)}
\emph{Meaning of the interval.} The credible interval $[-0.0528,-0.0358]$ states that, given the data and the prior, $A$ lies in it with $68\%$ probability; the confidence interval $[-0.0526,-0.0356]$ is the set of values that the data do not reject at the $32\%$ level, and the $68\%$ is the fraction of repeated experiments in which such an interval contains the true value. The numbers are the same here; the statements are not, and only the second is a property of the procedure that can be checked by repeating the experiment, which is what the bootstrap and toy checks did.

\subsection{Comparison: nuisance parameters and validation (complete)}
\emph{Nuisance parameters.} The frequentist analysis profiles the eleven nuisance parameters, the Bayesian one marginalizes them. The two differ when the likelihood is far from Gaussian or the prior is informative; here $A$ is almost uncorrelated with the others (largest correlation $-0.17$, with $A_b$), and the results coincide.

\emph{Validation.} The Bayesian toys are generated from the fitted model, which makes them a closure test; the bootstrap replicas of the frequentist check are resampled from the data and carry the features that the model does not describe. That both give the nominal coverage supports the interval, but they test different things (Sec.~\ref{sec:freq_coverage}).

\subsection{Conclusions: limitations (complete)}
Three limitations bound the result and indicate the next steps (Appendix~\ref{sec:limitations}): no tested model reaches $\chi^2/\mathrm{ndof}\approx1$, although $\Araw$ does not depend on the signal shape; no correction for the production and detection asymmetries is applied, so $\Araw\neq\ACP$; and the local Dalitz-plane asymmetries are counts in a region, without a background subtraction by event weights. The first of these, a model of the partially reconstructed and misidentified backgrounds, would also make the machine-learning selection usable in a fit.

\subsection{Model: label switching of the background components (complete)}
The mixture is invariant under the exchange $(\lambda_1,f)\leftrightarrow(\lambda_2,1-f)$, which would give the posterior two identical modes (label switching); the ordering $\lambda_1>\lambda_2$ removes this symmetry.

\subsection{Conclusions: lessons of each method, short version (complete)}
Each method taught something specific.
\begin{itemize}[leftmargin=*,itemsep=2pt]
\item \emph{Frequentist.} With the other eleven parameters profiled, Wilks' theorem gives a $68\%$ interval with the nominal coverage ($68.4\pm1.5\%$ in $1000$ bootstrap replicas, $0.70\pm0.01$ and $0.95\pm0.01$ in toy MC, where the estimator is unbiased and reaches the minimum variance bound); Hesse and Minos agree, so the likelihood is parabolic in $A$.
\item \emph{Bayesian.} With flat priors the posterior is the likelihood, and the credible interval coincides with the confidence interval and has the same coverage; the Bayesian analysis adds sampling diagnostics that must be checked (autocorrelation time, dispersed starts), a Bayes factor that depends on the prior width, and a posterior predictive check whose charge statistics, being fitted, have no power.
\item \emph{Machine learning (Appendix~\ref{sec:ml}).} Weak labels from the mass spectrum train classifiers of test AUC $\simeq0.86$; the selection leaves $\Araw$ unchanged, does not improve the precision where the background is small and cannot remove the misidentified decays under the signal.
\end{itemize}

\subsection{Conclusions: systematic study, short version (complete)}
The dominant systematic contribution is the particle-identification requirement, consistent with a statistical fluctuation of a nested sample; the background model is second and the signal shape does not matter, and the same variants and rule give the same uncertainty in the two methods. Several statistics depend on choices: the $p$-value of a binned $\chi^2$ on the binning, the Bayes factor on the prior, the coverage on whether the toy experiments come from the model or from the data.

\subsection{Data: charge combinations (complete)}
In this dataset only the charge combinations $(+1,-1,-1)$ and $(-1,+1,+1)$ occur for $(H1,H2,H3)$: $H1$ is always the kaon of charge opposite to the $B$, and $H2$, $H3$ are the interchangeable same-sign pair. The Dalitz-plot variables are therefore the invariant masses of the two opposite-sign pairs, $m^2(H1,H2)$ and $m^2(H1,H3)$.

\subsection{Comparison: machine-learning cross-check (complete)}
A classifier trained on weak labels from the mass spectrum (test AUC $0.86$) was used as an independent selection: a cut at $0.5$ keeps $92\%$ of the fitted signal and $36\%$ of the background and leaves $\Araw$ within its expected statistical spread ($-0.0451\pm0.0087$ Bayesian, $-0.0447\pm0.0086$ frequentist), without improving the precision; it cannot remove the misidentified decays that peak under the signal. Details, including the $K$-means partition of the Dalitz plane, are in Appendix~\ref{sec:ml}.

\subsection{Systematic study: list of the variants (complete)}
Each variant is fitted by both methods on exactly the same candidates.
\begin{itemize}[leftmargin=*,itemsep=1pt]
\item \emph{Window}: lower edge at $5050$ or $5150\,\MeV$, upper edge at $5450$ or $5700\,\MeV$.
\item \emph{Background}: a single exponential per charge; a double exponential with the shape shared between charges; a charge-symmetric background ($A_b=0$).
\item \emph{Signal}: a Crystal Ball shape with free tail parameters, Eq.~\eqref{eq:cb}; $\mu$ and $\sigma$ separate for the two charges.
\item \emph{Preselection}: one cut at a time loosened or tightened, $\texttt{ProbK}>0.4$ or $>0.6$, $\texttt{ProbPi}<0.4$ or $<0.6$, on all three tracks.
\end{itemize}
